\scanpage{236}
\mspage{109}
% archive sheet number: 235

\uline{Lemme} 8.6. Soit $\unsure{\mathcal{M}|M}$ une catégorie, $(C,\mathrm{TF})$
un précouple de Quillen dans $\mathcal{M}$, et $W\subset\mathrm{Fl}\,\mathcal{M}$ une partie de $\mathrm{Fl}\,\mathcal{M}$, stable par facteurs directs.
\deleted{Pour qu'on ait $\mathrm{TF}\subset W$, il faut et \unsure{s.} que toute flèche $f$ \unsure{a}} \inserted{Conditions équivalentes}

$a)$ \deleted{\unsure{tt}} $\mathrm{TF}\subset W$

$b)$ toute flèche $f$ dans $\mathcal{M}$ se factorise en $f=pi$, avec $i\in C$, $p\in W$.

$b')$ Comme $b)$, avec $f$ dans $\mathrm{TF}$.

Démonstration : les implications $a)\Rightarrow b)\Rightarrow b')$
sont tautologiques, et ne font pas appel à l'hypothèse \inserted{de stabilité} sur $\unsure{W_a|W.|W_o}$
Mais pour un $f\in\mathrm{TF}$ donné, pour voir qu'il est dans
$W$, \deleted{on utilise l'hyp} si $W$ est stable par facteurs directs, il suffit en effet de
le factoriser en $pi$, avec $i\in\unsure{\mathrm{TC}|C}$, $p\in W$.
En effet, \deleted{\unsure{on}} voit aussitôt qu'alors
$f$ est facteur direct de $p$ (car $i\in C$) et comme $p\in W$, on a $i\in W$, \unsure{qed.}

Revenons à 8.4. dém de $d)\Rightarrow a)$.

\deleted{On veut prouver que $\mathrm{TF}\subset W$, et on est ramené à prouver que toute flèche $f$ de $A^\wedge$ se factorise en $pi$, avec $i\in C$ i.e. $i$ mono, et $p\in W_A$}
